% ============================ INTRODUCTION GÉNÉRALE ==================
\chapter*{Introduction générale}
\addcontentsline{toc}{chapter}{Introduction générale}
\markboth{Introduction générale}{Introduction générale}

Les organisations d'ingénierie produisent et consomment des volumes considérables de
documentation technique : articles de recherche, manuels d'équipements, guides de normes,
schémas et spécifications. Cette connaissance, essentielle au travail quotidien des
ingénieurs, reste pourtant difficile d'accès : elle est enfermée dans des fichiers PDF ou
Word hétérogènes, dont une partie significative de l'information est portée par des
\emph{éléments visuels} (diagrammes d'architecture, schémas électriques, courbes de
résultats, tableaux) que les moteurs de recherche classiques, fondés sur le texte seul,
ignorent purement et simplement.

L'émergence des grands modèles de langue (LLM) a ouvert une voie nouvelle : le paradigme
RAG (\emph{Retrieval-Augmented Generation}) permet de coupler la capacité de synthèse d'un
LLM à une base documentaire propre à l'organisation, en récupérant à chaque question les
passages pertinents du corpus et en contraignant le modèle à y ancrer sa réponse. Le RAG
répond ainsi aux deux limites majeures des LLM utilisés seuls : l'absence de connaissance
des documents internes et la propension à l'hallucination. Toutefois, la majorité des
systèmes RAG restent \emph{unimodaux} : ils n'indexent que le texte, et une question dont la
réponse se trouve dans une figure reste sans réponse.

C'est précisément le problème auquel s'attaque ce projet, réalisé dans le cadre d'un stage
de Projet de Fin d'Année de huit semaines au sein d'\textbf{ALTEN}, leader mondial de
l'ingénierie et du conseil en technologies. L'objectif du stage est de concevoir et de
réaliser un \textbf{assistant RAG multimodal} capable d'interroger en langage naturel un
corpus de documents techniques et de restituer des réponses sourcées mêlant extraits de
texte et images pertinentes. Le projet couvre l'intégralité de la chaîne : ingestion et
analyse structurelle des documents (Docling), description automatique des figures par un
modèle de vision-langage, indexation vectorielle unifiée (ChromaDB), recherche hybride
dense + lexicale, génération de réponses citées (Mistral AI), exposition par une API REST
(FastAPI), interface de démonstration (React/TypeScript) et évaluation quantitative de la
qualité des réponses (RAGAS).

Une particularité de ce stage mérite d'être soulignée dès l'introduction : le développement
s'est déroulé sur un \emph{poste professionnel verrouillé}, derrière un proxy d'inspection
TLS (Zscaler), sans droits d'administration ni possibilité d'installer des services
système. Ces contraintes, représentatives des environnements industriels réels, ont
fortement structuré les choix techniques : préférence pour des API cloud accessibles,
solutions embarquées sans serveur (SQLite, ChromaDB persistant, FakeRedis), modèles
pré-téléchargés hors ligne ; ils constituent en eux-mêmes un retour d'expérience précieux.

Le présent rapport s'articule en trois chapitres. Le \textbf{chapitre 1} dresse l'état de
l'art des systèmes RAG et de leurs extensions multimodales, analyse les approches et outils
existants, et positionne la solution proposée. Le \textbf{chapitre 2} présente l'analyse
des besoins, la conception et l'architecture du système : pipeline fonctionnel,
architecture logicielle, modélisation UML, modèle de données et justification des choix
technologiques. Le \textbf{chapitre 3} détaille la réalisation : environnement de
développement contraint, implémentation des composants, intégration des services tiers,
sécurité applicative, évaluation RAGAS, résultats obtenus, difficultés rencontrées et
limites. Une conclusion générale et des perspectives d'évolution clôturent le document.
